Due importanti spazi metrici particolari sono gli spazi di Banach e gli spazi di Hilbert

\textbf{Spazi di Banach}

Uno spazio vettoriale normato si dice spazio di Banach se e solo se il relativo spazio metrico
definito con la distanza indotta dalla norma è completo.

Più schematicamente, sia $\left(\mathbb{X}, \|\cdot\|\right)$ uno spazio vettoriale normato,
definiamo la distanza in $\mathbb{X}$ indotta dalla norma come $d\left(x,y\right):=\|x-y\|$.
Allora $\left(\mathbb{X}, \|\cdot\|\right)$ si dice spazio di Banach se e solo se lo spazio
metrico $\left(\mathbb{X}, d\right)$ risulta completo.



\textbf{Spazi di Hilbert}

Uno spazio vettoriale dotato di prodotto scalare si dice spazio di Hilbert se e solo se
il relativo spazio metrico definito con la distanza indotta dalla norma, a sua volta
indotta dal prodotto scalare, è completo.

Più schematicamente, sia $\left(\mathbb{X}, \langle\cdot,\cdot\rangle\right)$ uno spazio
vettoriale dotato di prodotto scalare, definiamo la norma su $\mathbb{X}$ indotta dal prodotto
scalare come $\|x\|=\sqrt{\langle x,x\rangle}$ e, definita la norma, definiamo la distanza in
$\mathbb{X}$ indotta dalla norma come $d\left(x,y\right):=\|x-y\|$.
Allora $\left(\mathbb{X}, \langle\cdot,\cdot\rangle\right)$ si dice spazio di Hilbert se e solo
se lo spazio metrico $\left(\mathbb{X}, d\right)$ risulta completo.

\begin{observation}
    Uno spazio di Hilbert risulta a tutti gli effetti uno spazio anche di Banach, dal momento
    che é uno spazio di Banach dove la norma presa generalmente come data viene in realtà
    definita/indotta tramite il prodotto scalre dato (la distanza indotta poi dalla norma
    di uno spazio di Hilbert rispetta la medesima condizione richiesta alla distanza indotta
    dalla norma di uno spazio di Banach).
\end{observation}